\documentclass{article}	%<nh>
	%<*nhs>
\usepackage{graphicx}
\usepackage{graphpap}


\title{Hardware and Basic I/O}


	%<*hn>
\usepackage{fancyhdr}


%%%%%%%%%%%%%%%%%%%%%%%%%%%%%%%%%%%%%%%%5
%
%  Set Up Margins
\input{/home/blake/templates/pagedim.tex}


%%%%%%%%%%%%%%%%%%%%%%%%%%%%%%%%%%%%%%%%%%%%%%%%%
%
%         Page format Mods HERE
%
%Mod's to page size for this document
\addtolength\textwidth{0cm}
\addtolength\oddsidemargin{0cm}
\addtolength\headsep{0cm}
\addtolength\textheight{0cm}
%\linespread{0.894}   % 0.894 = 6 lines per inch, 1 = "single",  1.6 = "double"


%%%%%%%%%%%%%%%%%%%%%%%%%%%%% HEADER / FOOTER
\pagestyle{fancy}
%%%%%  Page header/footer fields for "NSF Style" proposal
%%%%%  \chead will be changed with each section
\lhead{\small\sc EE472}
\rhead{}
\chead{\large\bf Lecture Notes: Hardware and Basic I/O}
\lfoot{Hannaford, U. of Washington}
\rfoot{Oct. 2004}
\cfoot{\thepage}
%\renewcommand\headrulewidth{1pt}
%\renewcommand\footrulewidth{1pt}


\linespread{1.0}

	%<*>
\begin{document}



\section*{Paradigm C++ Problems and Procedures}
\section*{The Problem}




\section{Bsic Machine Organization}
\subsection{Registers}

[FIGURE - Registers]

\subsection{Busses}

A parallel bi-directional data path.

[FIGURE = big picture of bus]

Address: N bits which specify {\it which location}.

Data:   Contents to/from memory or I/O device

Control:  'traffic signals'

Bus Logic Notation:  [FIGURE]

Problem:  Multiple talkers

Solution:  Tri-state logic.

  [Figure: tri-state gate]

\subsection{Bus Timing}

Write Cycle

      [FIGURE]

Read Cycle

      [FIGURE]

\subsection{Microprocessor: Register View}

 [FIGURE:  PC, IR, Addr, GnlPrps, ALU]

 ALU = Arithmetic \& Logic Unit

\begin{itemize}

\item PC:  Program Counter. Holds Address of next instruction.
\item IR:  Instruction Register. Holds current instruction.
\item ADDR: Address Register. Holds address of next bus access.
\item GP\#: General Purpose Registers.  Hold intermediate results.
\item A1, A2: ALU Arguments. Hold inputs to an ALU operation.
\item RES:  Holds result of ALU operation.

\end{itemize}

\subsection{Steps to add two numbers}

\begin{enumerate}

\item C statement: {\tt c = a + b ;}.
\item  Load addr of {\tt a} into ADDR
\item  Wait for data from memory
\item  Clock data into A1
\item  Load addr of {\tt b} into ADDR
\item  Wait for data from memory
\item  Clock data into A2
\item  Send ADD command to ALU
\item  Load addr of {\tt c}
\item  Transfer RES to data bus
\item  Wait for data write to memory.

\end{enumerate}

\section{Machine Language Instructions}

The above example may be controlled by one or more machine instructions.
A machine instruction is binary data typically broken up into fields:

The Operation Code (Op-Code)

One, two, or three Operands.

Each instruction is typically between one and eight bytes.

Examples:           [FIGURE]

\subsection{Example: Add two numbers}
Note:  All assembler below is pseudo-code!

C: {\tt c = a + b;}

Assembler:
\begin{tabular}{|r|l|l|}
\hline
mem addr        & instruction 	& comment\\
\hline
0xA000          & MOV  a, A1   & move mem location a to ALU Arg 1 \\
\hline
0xA002          & MOV  b, A2    & move mem location b to ALU Arg 2 \\
\hline
0xA004          & ADD           & a one-byte instruction \\
\hline
0xA005          & MOV  r, c     & move ALU result to mem location c\\
\hline
\end{tabular}

There are many more types of operands in most architectures.  Examples:

\begin{itemize}
\item INDEXED
\begin{verbatim}
MOV   0xCA04(R2), G1         \\move data from memory location
                             \\ 0xCA04 + contents of R2 to register G1
\end{verbatim}

\item REGISTER INDIRECT
\begin{verbatim}
MOV   (R2), G1        	     \\move data from memory location
                             \\ contents of R2
\end{verbatim}

\end{itemize}


%%%%%%%%%%%%%%%%%%%%%%%%%%%%%%%%%%%%%%%%%%%%%%%%%%%%%%%%%%%%%%%%%%%%%%%%%%%%
\section{Basic I/O programming}

Intel I/O instructions

{\tt IN         R, char }
where {\tt char} is an 8-bit constant

{\tt OUT        char, R }

In intel architecture, {\tt R} can be {\tt AL, AX, or EAX} registers.



\end{document}

